\documentclass[journal,twocolumn]{IEEEtran}
\usepackage[utf8]{inputenc}
\usepackage[T5]{fontenc}
\usepackage[vietnamese,english]{babel}
\usepackage{amsmath,amssymb,amsfonts}
\usepackage{algorithmic}
\usepackage{graphicx}
\usepackage{textcomp}
\usepackage{booktabs}
\usepackage{hyperref}
\usepackage{multirow}
\usepackage{array}
\usepackage{float}

\begin{document}

\title{Class-wise Embedding-Guided Specimen Splitting for Leakage-Free Fine-Grained Wood Species Identification}

\author{Viet-Anh Le, and Collaborators}

\maketitle

\begin{abstract}
Nhận diện tự động loài gỗ là một nhiệm vụ quan trọng trong giám sát thương mại lâm sản, truy xuất nguồn gốc gỗ quý và hỗ trợ thực thi Công ước CITES trong bối cảnh khai thác gỗ bất hợp pháp vẫn là mối đe dọa lớn đối với đa dạng sinh học. Tuy nhiên, bài toán này đối mặt với hai thách thức cốt lõi. Thứ nhất, việc phân tích hình ảnh mặt cắt gỗ thuộc nhóm phân loại thị giác chi tiết, trong đó các loài cùng chi có cấu trúc mạch, tia gỗ và nhu mô rất tương đồng, trong khi biến thiên nội bộ loài lại lớn do tuổi cây, vùng sinh trưởng, vị trí cắt và điều kiện chụp. Thứ hai, các quy trình đánh giá sử dụng phương pháp phân chia ngẫu nhiên ở cấp ảnh thường gây rò rỉ dữ liệu cấp mẫu vật, do nhiều hình ảnh được cắt hoặc chụp từ cùng một mẫu vật lý bị phân tán đồng thời vào cả tập huấn luyện và tập kiểm thử. Điều này khiến mô hình ghi nhớ các đặc trưng lối tắt như vết xước, màu sắc, hướng cắt hoặc nhiễu cục bộ của mẫu thay vì học được các đặc trưng giải phẫu có khả năng tổng quát hóa.

Để giải quyết các hạn chế trên, nghiên cứu này đề xuất phương pháp phân chia mẫu vật dựa trên không gian nhúng theo từng loài (Class-wise Embedding-Guided Specimen Split - CEGS-Split), một giao thức chia dữ liệu nhằm cô lập mẫu vật lý và giảm thiểu rò rỉ ngữ nghĩa giữa tập huấn luyện và tập kiểm thử. Phương pháp CEGS-Split biểu diễn mỗi thư mục con như một nhóm mẫu vật vật lý độc lập, trích xuất đặc trưng nhúng bằng các mạng trích xuất đặc trưng cơ sở mạnh như EfficientNetV2 và Swin Transformer, sau đó chọn chiến lược phân chia dữ liệu phù hợp cho từng loài từ bốn chiến lược cốt lõi: phân chia dựa trên khoảng cách Mahalanobis, phân cụm phân cấp, phân nhóm có phân tầng và phân cụm phân tầng. Chúng tôi xây dựng thực nghiệm đối chiếu trên bộ dữ liệu S3 gồm 19 loài thuộc 6 chi, so sánh 14 hàm loss học khoảng cách và 4 thuật toán học tự giám sát trong cùng một quy trình đánh giá thống nhất. Giao thức đề xuất tạo ra một môi trường kiểm thử không rò rỉ dữ liệu, phản ánh chính xác hiệu năng thực tế, đồng thời đánh giá chất lượng không gian nhúng qua các chỉ số truy xuất và phân cụm như Recall@k, mAP, Silhouette, Davies-Bouldin, Dunn và NMI.
\end{abstract}

\begin{IEEEkeywords}
Fine-grained visual classification, Wood species identification, Metric learning, Data leakage, Deep learning, Self-supervised learning, CITES.
\end{IEEEkeywords}

\section{Introduction}

\subsection{Background}
Khai thác và buôn bán gỗ trái phép gây thiệt hại nghiêm trọng cho hệ sinh thái rừng, làm suy giảm đa dạng sinh học và tạo áp lực lên cSự phát triển của thị giác máy tính và học sâu mở ra khả năng tự động hóa nhận diện gỗ từ hình ảnh mặt cắt vĩ mô hoặc cận vi mô. Thay vì chỉ dựa vào mô tả thủ công của mạch gỗ, tia gỗ và nhu mô, mô hình học sâu có thể tự động học biểu diễn phân biệt trực tiếp từ dữ liệu hình ảnh. Tuy nhiên, đối với hình ảnh gỗ, một quy trình học sâu chỉ thực sự có ý nghĩa thực tế nếu nó đánh giá được khả năng nhận diện trên mẫu vật lý chưa từng xuất hiện trong quá trình huấn luyện, chứ không chỉ trên các lát cắt khác của cùng một khối gỗ.

\subsection{Problem Statement}
Nhận diện loài gỗ là bài toán phân loại thị giác chi tiết. Khoảng cách hình thái giữa các loài cùng chi thường rất nhỏ: các đặc trưng như phân bố lỗ mạch, dạng nhu mô quanh mạch, mật độ tia gỗ hoặc cấu trúc thớ có thể tương đồng đến mức khó phân biệt bằng mắt thường. Ngược lại, trong cùng một loài có thể xuất hiện biến thiên rất lớn do gỗ lõi và gỗ giác, độ tuổi cây, điều kiện thổ nhưỡng, độ ẩm, góc cắt và điều kiện chiếu sáng. Do đó, mô hình cần vừa học được đặc trưng ổn định của loài, vừa chịu được sự dịch chuyển phân phối ở cấp mẫu vật.

Về mặt dữ liệu, giả sử $\mathcal{X}=\{x_1,x_2,\dots,x_M\}$ là tập ảnh và $\mathcal{S}=\{S_1,S_2,\dots,S_N\}$ là tập mẫu vật lý, trong đó mỗi mẫu vật thường được lưu dưới dạng một thư mục con. Một ảnh $x_i$ có nhãn loài $y_i$ và thuộc về đúng một mẫu vật $S(x_i)\in\mathcal{S}$. Nếu chia ngẫu nhiên ở cấp ảnh, các ảnh
\begin{equation}
\{x_{j,1},x_{j,2},\dots,x_{j,m}\}\in S_j
\end{equation}
có thể xuất hiện đồng thời trong cả tập huấn luyện, tập kiểm định và tập kiểm thử. Khi đó, tập kiểm thử không còn độc lập về mặt mẫu vật, làm sai lệch nghiêm trọng kết quả đánh giá thực tế.

\subsection{Limitations of Prior Works}
Nhiều nghiên cứu nhận diện gỗ trước đây báo cáo độ chính xác rất cao khi áp dụng phương pháp phân chia ngẫu nhiên ở cấp ảnh. Tuy nhiên, trong các cơ sở dữ liệu gỗ thực tế, nhiều hình ảnh thường được chụp từ cùng một tiêu bản hoặc cắt ra từ cùng một khối vật lý. Mạng thần kinh tích chập hoặc Vision Transformer có năng lực ghi nhớ mạnh, nên khi một mẫu vật xuất hiện đồng thời trong cả tập huấn luyện và tập kiểm thử, mô hình có thể khai thác các lối tắt phi phân loại học như màu nền, vết xước, vết dao cắt, độ bóng bề mặt hoặc nhiễu camera. Các dấu hiệu này không phải đặc trưng phân loại học, vì vậy kết quả kiểm thử trong phòng thí nghiệm có thể bị thổi phồng và sẽ sụt giảm mạnh khi triển khai trên các khối gỗ hoàn toàn mới.

Một hạn chế khác là phần lớn các nghiên cứu hiện tại sử dụng bộ phân loại dựa trên hàm lỗi entropy chéo đơn thuần. Cách tiếp cận này chỉ tối ưu hóa biên quyết định khép kín cho các lớp đã biết, nhưng chưa trực tiếp tổ chức không gian nhúng đặc trưng một cách tối ưu. Với dữ liệu gỗ, việc biểu diễn không gian nhúng thông qua học khoảng cách sâu (deep metric learning) hoặc học tự giám sát đóng vai trò quan trọng để đảm bảo khả năng truy xuất tập mở (open-set retrieval).ặt mẫu vật, làm sai lệch kết quả đNghiên cứu này đặt mục tiêu xây dựng một giao thức nhận diện gỗ có khả năng tổng quát hóa trên mẫu vật mới. Trước hết, chúng tôi áp dụng nguyên tắc cô lập ở cấp mẫu vật (specimen-level isolation): toàn bộ hình ảnh thuộc cùng một thư mục con chỉ được phép xuất hiện trong một tập dữ liệu duy nhất. Điều kiện rời rạc được phát biểu như sau:
\begin{equation}
\mathcal{S}_{\text{train}}\cap\mathcal{S}_{\text{val}}=\emptyset,\quad
\mathcal{S}_{\text{train}}\cap\mathcal{S}_{\text{test}}=\emptyset,\quad
\mathcal{S}_{\text{val}}\cap\mathcal{S}_{\text{test}}=\emptyset.
\end{equation}
Từ đó, các tập dữ liệu thành phần được xây dựng độc lập:
\begin{equation}
\mathcal{D}_{\text{train}}=\{x_i\in\mathcal{X}\mid S(x_i)\in\mathcal{S}_{\text{train}}\},
\end{equation}
\begin{equation}
\mathcal{D}_{\text{val}}=\{x_i\in\mathcal{X}\mid S(x_i)\in\mathcal{S}_{\text{val}}\},\quad
\mathcal{D}_{\text{test}}=\{x_i\in\mathcal{X}\mid S(x_i)\in\mathcal{S}_{\text{test}}\}.
\end{equation}

Trên nền tảng đó, chúng tôi đề xuất giao thức Class-wise Embedding-Guided Specimen Split (CEGS-Split), một phương pháp lai tự động lựa chọn chiến lược phân chia dữ liệu phù hợp theo từng loài dựa trên phân tích không gian đặc trưng nhúng. Cách tiếp cận này chuyển đổi trọng tâm đánh giá từ cấp độ hình ảnh ngẫu nhiên sang cấp độ cô lập mẫu vật lý, giúp phản ánh trung thực năng lực của mô hình đối với các khối gỗ mới.

\subsection{Key Contributions}
Các đóng góp chính của nghiên cứu gồm:
\begin{itemize}
    \item Đóng góp bộ dữ liệu hình ảnh gỗ vĩ mô S3 được chuẩn hóa cấu trúc phân cấp (loài - mẫu vật - hình ảnh), tập trung vào 19 loài gỗ thuộc 6 chi có giá trị thương mại cao hoặc nằm trong danh mục cần bảo vệ nghiêm ngặt của CITES, hỗ trợ công tác nghiên cứu và giám định lâm sản.
    \item Đề xuất giao thức phân chia mẫu vật dựa trên không gian nhúng theo từng loài (CEGS-Split) nhằm loại bỏ triệt để hiện tượng rò rỉ dữ liệu cấp mẫu vật bằng cách sử dụng độ tương đồng đặc trưng để tối ưu hóa việc phân chia các thư mục con.
    \item Xây dựng hệ thống thực nghiệm đối chiếu toàn diện cho 14 hàm loss học khoảng cách sâu và 4 phương pháp học tự giám sát trên cùng một quy trình đánh giá cô lập mẫu vật.
    \item Đánh giá chất lượng mô hình đa chiều bằng cả chỉ số truy xuất hình ảnh và chỉ số hình học của không gian nhúng, bao gồm Recall@k, Precision@k, mAP, AUC, Silhouette, Davies-Bouldin, Calinski-Harabasz, Dunn và NMI.
\end{itemize}

\section{Related Work}

\subsection{Wood Species Identification}
Các phương pháp nhận diện gỗ truyền thống sử dụng đặc trưng thủ công như LBP, GLCM, histogram màu hoặc mô tả hình thái mạch gỗ. Giai đoạn gần đây, CNN như VGG, ResNet, EfficientNet, MobileNet và ConvNeXt được dùng rộng rãi để học đặc trưng trực tiếp từ ảnh. Các kiến trúc Transformer cũng cho thấy tiềm năng nhờ khả năng mô hình hóa quan hệ dài hạn trong texture. Dù vậy, nhiều nghiên cứu vẫn dùng random split ở cấp ảnh, chưa kiểm soát đầy đủ rò rỉ mẫu vật. Vì vậy, độ chính xác cao chưa chắc phản ánh khả năng nhận diện trên block gỗ mới.

\subsection{Deep Metric Learning}
Deep metric learning học một ánh xạ $f_\theta(x)$ sao cho các mẫu cùng lớp gần nhau và các mẫu khác lớp xa nhau trong không gian nhúng. Các loss cổ điển như Contrastive Loss và Triplet Loss tối ưu trực tiếp khoảng cách Euclid, nhưng thường phụ thuộc vào mining và có thể hội tụ chậm khi batch chứa nhiều cặp dễ. Các phương pháp hiện đại như Multi-Similarity, Circle Loss, ArcFace, SubCenter ArcFace, Proxy Anchor và SoftTriple đưa thêm trọng số động, margin góc hoặc proxy học được để làm biên phân tách sắc nét hơn. Đối với ảnh gỗ, DML đặc biệt phù hợp vì mục tiêu cuối cùng không chỉ là gán nhãn mà còn là truy vấn mẫu tương tự và kiểm tra cấu trúc cụm của các loài.

\subsection{Self-Supervised Representation Learning}
Self-supervised learning học biểu diễn từ các biến đổi của ảnh mà không phụ thuộc hoàn toàn vào nhãn thủ công. SimCLR tối ưu InfoNCE bằng mẫu âm, trong khi BYOL và SimSiam dùng predictor và stop-gradient để tránh collapse mà không cần negative pairs. Barlow Twins tối ưu ma trận tương quan chéo để vừa kéo hai view của cùng ảnh lại gần, vừa giảm dư thừa giữa các chiều đặc trưng. Với ảnh gỗ, SSL có khả năng học texture, hướng thớ và cấu trúc mạch bền vững trước biến đổi ảnh; tuy nhiên cần được đánh giá trong bối cảnh không rò rỉ cấp mẫu vật.

\section{Proposed Methodology}

\subsection{Overall Framework}
Quy trình nghiên cứu đề xuất gồm sáu bước chính. Thứ nhất, dữ liệu S3 được chuẩn hóa hình ảnh và loại bỏ các mẫu không có nhãn định danh cụ thể ở cấp loài. Thứ hai, các hình ảnh được nhóm theo thư mục con (subfolder) đại diện cho mẫu vật lý. Thứ ba, các đặc trưng biểu diễn ban đầu được trích xuất bằng mạng trích xuất đặc trưng cơ sở (backbone) mạnh như EfficientNetV2-M hoặc Swin-Large để lập bản đồ phân bố ngữ nghĩa của từng mẫu vật. Thứ tư, giao thức CEGS-Split tự động chọn phương pháp phân chia dữ liệu tối ưu cho từng loài gỗ nhằm thiết lập các tập huấn luyện (train), kiểm định (validation) và kiểm thử (test) cô lập mẫu vật lý, đảm bảo chống rò rỉ dữ liệu. Thứ năm, mạng ConvNeXt-Tiny được tinh chỉnh (fine-tune) kết hợp với bộ đầu chiếu (projection head) phi tuyến để học biểu diễn đặc trưng nhúng 256 chiều. Thứ sáu, mô hình được đánh giá toàn diện bằng các phép truy xuất hình ảnh (retrieval), thuật toán phân cụm (clustering), trực quan hóa giảm chiều t-SNE và bản đồ kích hoạt chú ý Grad-CAM/Finer-CAM.

\subsection{Phân chia mẫu vật dựa trên không gian nhúng theo từng loài (CEGS-Split)}
Với mỗi thư mục con $S_j$, đặc trưng nhúng đại diện được tính bằng trọng tâm (centroid) của các hình ảnh thuộc thư mục con đó:
\begin{equation}
E_j=\frac{1}{|S_j|}\sum_{x_i\in S_j} e_i,
\end{equation}
trong đó $e_i\in\mathbb{R}^{D}$ là vector đặc trưng của hình ảnh $x_i$. Nếu số chiều quá lớn so với số lượng thư mục con, thuật toán phân tích thành phần chính (PCA) được áp dụng để giảm chiều về $d'=\min(G-2,128)$ nhằm ổn định việc ước lượng ma trận hiệp phương sai và giảm thiểu hiện tượng quá khớp do số chiều lớn (curse of dimensionality).

Giao thức CEGS-Split hoạt động bằng cách quét qua một không gian các chiến lược phân chia tiềm năng (gồm các nhóm phương pháp chính: khoảng cách Mahalanobis, phân cụm phân cấp, phân nhóm có phân tầng và phân cụm phân tầng) để xác định thuật toán phù hợp nhất cho từng loài gỗ cụ thể. Trong đó, bốn chiến lược phân chia cốt lõi được phát triển nhằm duy trì tính cô lập thư mục con nghiêm ngặt ở các mức độ hình học và phân phối khác nhau. Phép phân chia ngẫu nhiên có phân tầng đóng vai trò phương pháp đối chiếu cơ sở (baseline) để so sánh hiệu năng, còn các phương pháp còn lại đều cố gắng giữ tính cô lập mẫu vật lý ở các mức độ khác nhau nhằm triệt tiêu rò rỉ thông tin. Phương pháp CEGS-Split tự động chọn phương pháp phân chia tối ưu theo từng loài thay vì áp dụng một quy tắc cứng nhắc duy nhất cho toàn bộ bộ dữ liệu.

\subsubsection{Phân chia dựa trên khoảng cách Mahalanobis}
Chiến lược này xem các mẫu vật ở xa trọng tâm (centroid) của lớp là các mẫu dị biệt (outlier) về mặt ngữ nghĩa và đưa chúng vào tập kiểm thử hoặc kiểm định nhằm thử nghiệm giới hạn (stress-test) khả năng tổng quát hóa của mô hình. Với tập mẫu chưa phân bổ $U$, centroid $\mu_U$ và hiệp phương sai $\Sigma_U$ được xác định bởi:
\begin{equation}
\mu_U=\frac{1}{|U|}\sum_{j\in U} z_j,
\end{equation}
\begin{equation}
\Sigma_U=\frac{1}{|U|-1}\sum_{j\in U}(z_j-\mu_U)(z_j-\mu_U)^T+\epsilon I.
\end{equation}
Khoảng cách Mahalanobis của mẫu $j$ đến centroid hiện tại là:
\begin{equation}
d_{\text{Mahal}}(z_j,\mu_U)=\sqrt{(z_j-\mu_U)^T\Sigma_U^{-1}(z_j-\mu_U)}.
\end{equation}
Tại mỗi vòng lặp, mẫu ở xa nhất $j^*=\arg\max_{j\in U}d_{\text{Mahal}}(z_j,\mu_U)$ được đưa vào tập kiểm thử hoặc kiểm định dựa trên kích thước yêu cầu, các mẫu còn lại được giữ lại làm dữ liệu huấn luyện.

\subsubsection{Phân chia dựa trên phân cụm phân cấp}
Chiến lược này phát hiện các nhóm mẫu vật tự nhiên trong cùng một loài bằng thuật toán phân cụm phân cấp liên kết tích tụ (Agglomerative Hierarchical Clustering) với tiêu chí liên kết Ward. Số lượng cụm tối ưu $K^*$ được xác định bằng phương pháp khuỷu tay (elbow method) áp dụng trên tổng bình phương khoảng cách trong cụm (Within-Cluster Sum of Squares - WCSS):
\begin{equation}
\text{WCSS}(K)=\sum_{k=1}^{K}\sum_{z_j\in\mathcal{C}_k}\|z_j-\mu_{\mathcal{C}_k}\|_2^2.
\end{equation}
After scaling $K$ and $\text{WCSS}(K)$ to $[0,1]$, the elbow point is selected by:
\begin{equation}
K^*=\arg\max_K \frac{|x_K+y_K-1|}{\sqrt{2}}.
\end{equation}
Mỗi cụm $\mathcal{C}_k$ được giữ nguyên vẹn khi phân bổ vào tập huấn luyện, kiểm định hoặc kiểm thử. Điều này ngăn các nhóm vân gỗ có cấu trúc bề mặt gần nhau bị chia cắt sang các tập khác nhau, giảm thiểu nguy cơ rò rỉ dữ liệu ở cấp độ cụm (cluster-level leakage).

\subsubsection{Phân chia dựa trên nhóm có phân tầng}
Phương pháp phân chia nhóm có phân tầng (Stratified Group Split) tối ưu hóa đồng thời hai ràng buộc: cô lập nhóm mẫu vật và cân bằng phân bố lớp giữa các tập. Gọi $w_{g,c}$ là số lượng ảnh của loài $c$ trong nhóm $g$, số lượng mẫu mục tiêu cho phân chia $f$ là:
\begin{equation}
T_{f,c}=\frac{1}{K}\sum_{g=1}^{G}w_{g,c}.
\end{equation}
Thuật toán gán tham lam sẽ lựa chọn tập:
\begin{equation}
f^*=\arg\min_f \sum_{c=1}^{C}\left(\frac{C_{f,c}+w_{g,c}}{T_{f,c}}-1\right)^2,
\end{equation}
trong đó $C_{f,c}$ là số lượng mẫu hiện tại của loài $c$ trong tập $f$. Phương pháp phân chia này phù hợp với các loài gỗ cần ưu tiên sự cân bằng số lượng mẫu giữa các tập để ổn định quá trình tính toán đạo hàm (gradient).

\subsubsection{Phân chia phân cụm phân tầng}
Phương pháp này kết hợp giữa phân cụm đặc trưng và phân chia phân tầng theo khoảng cách hình học. Sau khi gom cụm mẫu vật, khoảng cách từ trọng tâm cụm $\mu_{\mathcal{C}_k}$ đến trọng tâm chung của toàn loài $\mu_{\text{global}}$ được tính bằng:
\begin{equation}
d(\mathcal{C}_k) = \sqrt{ D^2(\mu_{\mathcal{C}_k}, \mu_{\text{global}}) },
\end{equation}
where:
\begin{equation}
D^2(\mu_{\mathcal{C}_k}, \mu_{\text{global}}) = (\mu_{\mathcal{C}_k}-\mu_{\text{global}})^T\Sigma_{\text{global}}^{-1}(\mu_{\mathcal{C}_k}-\mu_{\text{global}}).
\end{equation}
Các cụm được chia thành ba dải Near, Mid và Far. Khi phân bổ, thuật toán chọn tập có độ hụt lớn nhất so với mục tiêu:
\begin{equation}
\Delta_s=\max(0,N_s^{\text{target}}-N_s^{\text{current}}),\quad s\in\{\text{train},\text{val},\text{test}\}.
\end{equation}
Nhờ đó, cả ba tập huấn luyện, kiểm định và kiểm thử đều chứa đầy đủ các mẫu điển hình, mẫu trung gian và mẫu dị biệt, tránh việc dồn tất cả các mẫu ngoại lai (outlier) vào tập kiểm thử.

\subsection{CEGS-Split Configuration}
CEGS-Split là một cấu hình lai theo từng loài. Thay vì áp dụng một chiến lược duy nhất cho toàn bộ dataset, mỗi loài được gán phương pháp chia phù hợp với đặc điểm sinh học, số lượng subfolder và kết quả phân tích phân bố không gian biểu diễn. Cấu hình này được suy ra từ quy trình tìm kiếm tự động quét các chiến lược chia dữ liệu trên từng lớp, dùng EfficientNetV2-M để tạo embedding phục vụ thuật toán phân hoạch, dùng Swin-Large để đo nearest-neighbor similarity giữa train và test, rồi chọn phương pháp có mean nearest-neighbor similarity nhỏ nhất nhằm giảm thiểu tối đa nguy cơ rò rỉ ngữ nghĩa. Bảng \ref{tab:cegs_split} trình bày cấu hình chi tiết đang được triển khai thực tế.

\begin{table}[htbp]
\centering
\caption{Class-specific configuration of the CEGS-Split protocol}
\label{tab:cegs_split}
\resizebox{\columnwidth}{!}{%
\begin{tabular}{ll}
\toprule
\textbf{Species} & \textbf{Split Strategy} \\
\midrule
\textit{Afzelia africana} & Stratified Group \\
\textit{Afzelia bella} & Hierarchical Clustering \\
\textit{Afzelia pachyloba} & Agglomerative Stratified \\
\textit{Afzelia quanzensis} & Mahalanobis Iterative \\
\textit{Dalbergia cochinchinensis} & Agglomerative Stratified \\
\textit{Dalbergia melanoxylon} & Mahalanobis Iterative \\
\textit{Dalbergia oliveri} & Stratified Group \\
\textit{Dalbergia rimosa} & Hierarchical Clustering \\
\textit{Dalbergia tonkinensis} & Hierarchical Clustering \\
\textit{Guibourtia arnoldiana} & Hierarchical Clustering \\
\textit{Guibourtia coleosperma} & Agglomerative Stratified \\
\textit{Guibourtia ehie} & Hierarchical Clustering \\
\textit{Peltogyne pubescens} & Mahalanobis Iterative \\
\textit{Pterocarpus erinaceus} & Agglomerative Stratified \\
\textit{Pterocarpus indicus} & Agglomerative Stratified \\
\textit{Pterocarpus macrocarpus} & Hierarchical Clustering \\
\textit{Pterocarpus soyauxii} & Hierarchical Clustering \\
\textit{Sindora cochinchinensis} & Mahalanobis Iterative \\
\textit{Sindora tonkinensis} & Agglomerative Stratified \\
\bottomrule
\end{tabular}%
}
\end{table}


Đối với các loài gỗ có phân bổ hoán đổi (Swapped (Val)), kết quả chia tạm thời được điều chỉnh như sau:
\begin{equation}
\mathcal{D}_{\text{train}}=\mathcal{D}_{\text{train\_temp}},\quad
\mathcal{D}_{\text{val}}=\mathcal{D}_{\text{test\_temp}},\quad
\mathcal{D}_{\text{test}}=\mathcal{D}_{\text{val\_temp}}.
\end{equation}
Với các loài có cấu hình tiêu chuẩn (Standard (Test)), kết quả ban đầu được giữ nguyên. Cơ chế này giúp cân bằng độ khó giữa tập kiểm định và tập kiểm thử theo đặc thù của từng loài gỗ, đồng thời vẫn duy trì sự cô lập thư mục con nghiêm ngặt.

\subsection{Mạng trích xuất đặc trưng cơ sở và Bộ đầu chiếu phi tuyến}
Mô hình nhận diện chính sử dụng kiến trúc ConvNeXt-Tiny làm mạng trích xuất đặc trưng cơ sở (backbone). Khoảng $90\%$ các lớp đầu tiên được đóng băng (freeze) để giữ lại các đặc trưng tổng quát học từ tập ImageNet, tránh hiện tượng quá khớp (overfitting) trên các loài gỗ có số lượng mẫu vật hạn chế. Đặc trưng thô $\mathbf{f}\in\mathbb{R}^{768}$ sau đó được đưa qua bộ đầu chiếu (projection head) phi tuyến để tạo đặc trưng nhúng $\mathbf{z}\in\mathbb{R}^{256}$. Đối với bộ đầu chiếu gồm hai tầng liên kết đầy đủ, biểu diễn chuẩn hóa $L_2$ được tính toán như sau:
\begin{equation}
\tilde{\mathbf{z}}=W_2\text{ReLU}(W_1\mathbf{f}+b_1)+b_2,
\end{equation}
\begin{equation}
\mathbf{z}=\frac{\tilde{\mathbf{z}}}{\|\tilde{\mathbf{z}}\|_2}.
\end{equation}
Do $\|\mathbf{z}\|_2=1$, khoảng cách Euclid giữa hai đặc trưng nhúng liên hệ trực tiếp với độ tương đồng cosine (cosine similarity):
\begin{equation}
d(\mathbf{z}_a,\mathbf{z}_b)=\|\mathbf{z}_a-\mathbf{z}_b\|_2=\sqrt{2(1-\mathbf{z}_a^T\mathbf{z}_b)}.
\end{equation}

\subsection{Các phương pháp đối chiếu và công thức hàm loss của học khoảng cách}
Để đánh giá khách quan các hướng học biểu diễn đặc trưng, chúng tôi thực nghiệm đối chiếu 14 hàm loss học khoảng cách trong cùng một quy trình huấn luyện thống nhất về mạng cơ sở, bộ đầu chiếu, tăng cường dữ liệu và giao thức CEGS-Split. Ký hiệu chung: $\mathbf{z}_i\in\mathbb{R}^{D}$ là đặc trưng nhúng đã chuẩn hóa $L_2$, $y_i$ là nhãn loài, $d_{ij}=\|\mathbf{z}_i-\mathbf{z}_j\|_2$ là khoảng cách Euclid, $S_{ij}=\mathbf{z}_i^T\mathbf{z}_j$ là độ tương đồng cosine và $[u]_+=\max(0,u)$. Đối với hình ảnh mặt cắt gỗ, các hàm loss này không chỉ được đánh giá bằng độ chính xác (accuracy) mà còn bằng khả năng xây dựng không gian nhúng có cấu trúc phân cụm ổn định cho các loài gỗ trong cùng chi.

\subsubsection{Hàm loss Triplet cơ bản}
Hàm loss Triplet (Triplet Loss) tối ưu hóa mối quan hệ giữa bộ ba mẫu gồm mẫu neo (anchor) $a$, mẫu tích cực (positive) $p$ thuộc cùng loài và mẫu tiêu cực (negative) $n$ thuộc loài khác. Mục tiêu của hàm loss là thu hẹp khoảng cách giữa mẫu neo và mẫu tích cực, đồng thời đẩy mẫu tiêu cực ra xa hơn mẫu neo ít nhất một khoảng biên (margin) $\alpha$:
\begin{equation}
L_{\text{triplet}}=\frac{1}{|\mathcal{T}|}\sum_{(a,p,n)\in\mathcal{T}}\left[d(a,p)^2-d(a,n)^2+\alpha\right]_+.
\end{equation}
Đây là mô hình đối chiếu nền tảng để kiểm tra liệu mô hình có học được cấu trúc kéo-gần/đẩy-xa cơ bản hay không. Hạn chế của phương pháp này là số lượng bộ ba hợp lệ tăng nhanh theo kích thước lô dữ liệu $O(B^3)$ và phần lớn các bộ ba trong lô dữ liệu (batch) nhanh chóng trở thành bộ ba dễ (có giá trị loss bằng 0). Trên bộ dữ liệu S3, điều này đặc biệt bất lợi do các loài gỗ cùng chi tạo ra các ranh giới phân tách rất khó, trong khi số lượng mẫu trong một lô dữ liệu có thể không đủ đa dạng để tìm được các bộ ba hữu ích.

\subsubsection{Hàm loss Triplet với khai thác mẫu khó trung bình}
Phương pháp Triplet bán khó (Semi-hard Triplet Loss) lựa chọn các mẫu tiêu cực (negative) không quá dễ nhưng cũng không quá nhiễu. Với một cặp mẫu neo và mẫu tích cực $(a,p)$, mẫu tiêu cực $n$ được chọn nếu thỏa mãn:
\begin{equation}
d(a,p)^2<d(a,n)^2<d(a,p)^2+\alpha.
\end{equation}
Loss vẫn giữ dạng triplet chuẩn:
\begin{equation}
L_{\text{semi-hard}}=\frac{1}{|\mathcal{T}_{sh}|}\sum_{(a,p,n)\in\mathcal{T}_{sh}}\left[d(a,p)^2-d(a,n)^2+\alpha\right]_+.
\end{equation}
Cơ chế này giúp ổn định đạo hàm (gradient) hơn so với việc khai thác mẫu khó nhất (hardest mining), do mẫu tiêu cực vẫn nằm ngoài khoảng cách mẫu tích cực một chút, tránh cho mô hình bị ảnh hưởng bởi các mẫu dị biệt (outlier) hoặc nhãn bị nhiễu. Đối với hình ảnh mặt cắt gỗ, cơ chế khai thác mẫu bán khó (semi-hard mining) rất hữu ích khi cần học biên giới phân biệt giữa các loài gỗ gần nhau trong cùng chi \textit{Dalbergia} mà không làm mô hình bị tối ưu quá mức đối với các mẫu bất thường.

\subsubsection{Hàm loss Triplet biên mềm}
Phương pháp Triplet biên mềm (Soft-Margin Triplet) thay thế hàm hinge cứng bằng hàm log-sum-exp mịn:
\begin{equation}
L_{\text{soft-triplet}}=\frac{1}{N}\sum_{i=1}^{N}\log\left[1+e^{d(a_i,p_i)^2-d(a_i,n_i)^2}\right].
\end{equation}
Khác với hàm loss Triplet truyền thống có biên cố định, hàm này vẫn duy trì một lượng cập nhật đạo hàm nhỏ ngay cả khi mẫu tiêu cực đã ở khoảng cách tương đối an toàn. Điều này làm giảm độ nhạy của mô hình đối với việc chọn siêu tham số biên $\alpha$, giúp quá trình tối ưu hội tụ mượt mà hơn. Trong bối cảnh bộ dữ liệu S3, biên mềm đóng vai trò cầu nối hiệu quả giữa phương pháp đối chiếu Triplet cơ bản và các hàm loss hiện đại, giúp mô hình liên tục tinh chỉnh cấu trúc của không gian nhúng thay vì dừng học sớm khi đã vượt qua biên cứng.

\subsubsection{Hàm loss Angular}
Hàm loss Angular (Angular Loss) tối ưu hóa quan hệ góc của tam giác được tạo bởi bộ ba mẫu (triplet) thay vì chỉ tối ưu hóa khoảng cách Euclid đơn thuần. Với bộ ba mẫu $(a,p,n)$, hàm loss được xác định như sau:
\begin{equation}
L_{\text{angular}}=\frac{1}{|\mathcal{T}|}\sum_{(a,p,n)\in\mathcal{T}}\log\left[1+e^{\Delta_{\theta}(a,p,n)}\right],
\end{equation}
where:
\begin{equation}
\begin{aligned}
\Delta_{\theta}(a,p,n) = &4\tan^2\alpha(\mathbf{z}_a+\mathbf{z}_p)^T\mathbf{z}_n \\
&- 2(1+\tan^2\alpha)\mathbf{z}_a^T\mathbf{z}_p.
\end{aligned}
\end{equation}
Do đặc trưng nhúng đã được chuẩn hóa L2, biên góc ít bị ảnh hưởng bởi độ lớn (scale) của vector đặc trưng. Điều này rất phù hợp với hình ảnh gỗ vĩ mô, nơi cùng một cấu trúc thớ gỗ có thể thay đổi cường độ sáng, độ phóng đại hoặc độ ẩm bề mặt, nhưng quan hệ góc trong không gian nhúng vẫn bảo toàn được thông tin phân biệt loài.

\subsubsection{Hàm loss Multi-Similarity}
Hàm loss tương đồng đa chiều (Multi-Similarity Loss) khai thác đồng thời thông tin từ nhiều cặp tích cực và tiêu cực trong cùng một lô dữ liệu, đồng thời gán trọng số lớn hơn cho các cặp có độ khó cao:
\begin{equation}
L_{\text{MS}}=\frac{1}{N}\sum_{i=1}^{N} \left[ \phi^+(P_i) + \phi^-(N_i) \right],
\end{equation}
where:
\begin{equation}
\begin{aligned}
\phi^+(P_i) &= \frac{1}{\alpha}\log\left[1+\sum_{p\in P_i}e^{-\alpha(S_{ip}-\lambda)}\right], \\
\phi^-(N_i) &= \frac{1}{\beta}\log\left[1+\sum_{n\in N_i}e^{\beta(S_{in}-\lambda)}\right].
\end{aligned}
\end{equation}
Thành phần tích cực phạt các mẫu cùng loài chưa đủ gần nhau trong không gian đặc trưng; thành phần tiêu cực phạt các mẫu khác loài nhưng lại có độ tương đồng cosine lớn. Đây là một trong những hàm loss đối chiếu quan trọng nhất đối với bộ dữ liệu S3, do các loài gỗ khác nhau nhưng cùng chi thường tạo ra các mẫu tiêu cực khó (hard negative) tự nhiên trong thực tế, ví dụ giữa loài \textit{Dalbergia oliveri} và \textit{Dalbergia tonkinensis}.

\subsubsection{Hàm loss Contrastive cơ bản}
Hàm loss tương phản (Contrastive Loss) tối ưu hóa khoảng cách trên từng cặp mẫu riêng lẻ. Gọi $q_{ij}=0$ nếu hai mẫu $i, j$ cùng loài ($y_i=y_j$) và $q_{ij}=1$ nếu khác loài ($y_i\ne y_j$), hàm loss được phát biểu:
\begin{equation}
L_{\text{contrastive}}=(1-q_{ij})\frac{1}{2}d_{ij}^{2}+q_{ij}\frac{1}{2}[m-d_{ij}]_+^2.
\end{equation}
Hàm loss kéo các cặp hình ảnh cùng loài về gần nhau và đẩy các cặp khác loài ra ngoài khoảng biên $m$. Ưu điểm của phương pháp này là trực quan và dễ cài đặt; tuy nhiên, hạn chế lớn là một biên khoảng cách duy nhất được áp dụng chung cho mọi cặp mẫu tiêu cực, không phản ánh được mức độ khó phân biệt khác nhau giữa các loài cùng chi so với các loài thuộc các chi khác nhau.

\subsubsection{Hàm loss Contrastive với khai thác cặp khó trực tuyến}
Phương pháp tương phản khai thác mẫu khó trực tuyến (Online Hard Contrastive) tập trung tối ưu hóa trên cặp tích cực xa nhất và cặp tiêu cực gần nhất trong cùng một lô dữ liệu (batch):
\begin{equation}
d_{\text{pos\_max}}^i=\max_{p:y_p=y_i,p\ne i}d_{ip},\quad d_{\text{neg\_min}}^i=\min_{n:y_n\ne y_i}d_{in}.
\end{equation}
\begin{equation}
L_{\text{hard-cont}}=\frac{1}{B}\sum_{i=1}^{B}\left[(d_{\text{pos\_max}}^i)^2+[m-d_{\text{neg\_min}}^i]_+^2\right].
\end{equation}
Cách tiếp cận này tập trung tài nguyên tính toán đạo hàm vào các vùng ranh giới khó phân tách nhất thay vì các cặp mẫu dễ. Đối với hình ảnh gỗ, nó giúp mô hình tập trung phân biệt các mẫu gỗ có vân tương đồng giữa hai loài khác nhau, nhưng cũng nhạy cảm hơn với các mẫu dị biệt (outlier), ảnh bị nhiễu hoặc nhãn bị gán sai.

\subsubsection{Hàm loss cấu trúc Lifted}
Hàm loss cấu trúc Lifted (Lifted Structured Loss) mở rộng cơ chế học tương phản bằng cách đồng thời xem xét tất cả các mẫu tiêu cực liên quan đến một cặp mẫu tích cực $(i,j)$:
\begin{equation}
L_{\text{lifted}}=\frac{1}{2|P|}\sum_{(i,j)\in P}\left[d_{ij}+\log\left(\Phi(i, j)\right)\right]_+^2,
\end{equation}
where:
\begin{equation}
\begin{aligned}
\Phi(i, j) = &\sum_{k:y_k\ne y_i}e^{\alpha-d_{ik}} \\
&+ \sum_{k:y_k\ne y_j}e^{\alpha-d_{jk}}.
\end{aligned}
\end{equation}
Hàm $\log\sum\exp$ đóng vai trò là một xấp xỉ trơn của phép toán lấy giá trị cực đại, do đó các mẫu tiêu cực gần mẫu neo nhất sẽ tự động nhận lực đẩy mạnh nhất. Hàm loss này rất hữu ích để đánh giá hiệu quả của việc tối ưu hóa cấu trúc lô dữ liệu đối với bề mặt gỗ, nơi nhiều loài khác nhau thường có một số mẫu đặc trưng nằm rất gần nhau trong không gian nhúng.

\subsubsection{Hàm loss Circle}
Hàm loss Circle (Circle Loss) đề xuất một cách tiếp cận tối ưu hóa độ tương đồng thống nhất, trong đó các mẫu tích cực và tiêu cực được gán các trọng số tự điều chỉnh:
\begin{equation}
L_{\text{circle}} = \log \left[ 1 + \Omega_{\text{circle}}(i) \right],
\end{equation}
where:
\begin{equation}
\begin{aligned}
\Omega_{\text{circle}}(i) = &\sum_{n\in N_i}e^{\gamma\alpha_n(S_{in}-\Delta_n)} \times \\
&\sum_{p\in P_i}e^{-\gamma\alpha_p(S_{ip}-\Delta_p)},
\end{aligned}
\end{equation}
and:
\begin{equation}
\alpha_p=[O_p-S_{ip}]_+, \quad \alpha_n=[S_{in}-O_n]_+.
\end{equation}
Mẫu tích cực đã ở đủ gần sẽ nhận trọng số cập nhật nhỏ, trong khi mẫu tiêu cực quá gần sẽ bị đẩy đi rất mạnh. Cơ chế này đặc biệt phù hợp với bài toán phân loại chi tiết (fine-grained), do nó tránh việc tối ưu hóa quá mức trên các loài dễ phân biệt, đồng thời tập trung đạo hàm vào các cặp mẫu thuộc cùng một chi có độ tương đồng cao.

\subsubsection{Hàm loss Proxy Anchor}
Phương pháp Proxy Anchor thay thế việc so sánh trực tiếp giữa các hình ảnh bằng việc so sánh hình ảnh với các mẫu đại diện (proxy) của từng lớp. Với proxy $p$ đại diện cho một lớp, hàm loss được tính như sau:
\begin{equation}
L_{\text{PA}} = \frac{1}{|P^+|}\sum_{p\in P^+}\log \Gamma^+(p) + \frac{1}{|P|}\sum_{p\in P}\log \Gamma^-(p),
\end{equation}
where:
\begin{equation}
\begin{aligned}
\Gamma^+(p) &= 1+\sum_{\mathbf{z}\in X_p^+}e^{-\alpha(S(\mathbf{z},p)-\delta)}, \\
\Gamma^-(p) &= 1+\sum_{\mathbf{z}\in X_p^-}e^{\alpha(S(\mathbf{z},p)+\delta)}.
\end{aligned}
\end{equation}
Do số lượng lớp của bộ dữ liệu S3 tương đối nhỏ, phương pháp học dựa trên proxy (proxy-based learning) mang lại chi phí tính toán thấp và dòng đạo hàm ổn định. Ngoài ra, Proxy Anchor giúp giảm thiểu sự phụ thuộc vào cấu trúc thành phần của lô dữ liệu (batch composition) - một ưu điểm quan trọng khi các loài gỗ trong cơ sở dữ liệu có sự mất cân bằng lớn về số lượng ảnh và số lượng thư mục con.

\subsubsection{Hàm loss ArcFace}
Phương pháp ArcFace (Additive Angular Margin Loss) áp dụng trực tiếp một biên góc cộng (additive angular margin) vào giá trị logit của lớp chính xác:
\begin{equation}
\begin{aligned}
L_{\text{ArcFace}} = -\frac{1}{N}\sum_{i=1}^{N}\log \frac{e^{s\cos(\theta_{y_i}+m)}}{\Lambda_i},
\end{aligned}
\end{equation}
where:
\begin{equation}
\begin{aligned}
\Lambda_i = e^{s\cos(\theta_{y_i}+m)} + \sum_{j\ne y_i}e^{s\cos\theta_j}.
\end{aligned}
\end{equation}
Do cả đặc trưng nhúng và ma trận trọng số của lớp đều được chuẩn hóa, ArcFace thiết lập một biên khoảng cách trắc địa (geodesic margin) trên bề mặt siêu cầu đơn vị. Hàm loss này tỏ ra đặc biệt hiệu quả trong bài toán phân loại chi tiết (fine-grained), do nó ép cụm đặc trưng của cùng một lớp co lại trong một vùng góc hẹp và đẩy cụm của các lớp khác ra xa, rất thích hợp với các loài gỗ có sự khác biệt hình thái cực kỳ nhỏ.

\subsubsection{Hàm loss SubCenter ArcFace}
Phương pháp SubCenter ArcFace mở rộng thuật toán ArcFace bằng cách gán nhiều trọng tâm phụ (sub-center) cho mỗi lớp. Với cấu hình $K$ trọng tâm phụ cho mỗi loài, giá trị logit của lớp chính xác được tính toán dựa trên trọng tâm phụ gần nhất:
\begin{equation}
\cos\theta_{y_i}=\max_{k=1}^{K}\frac{\mathbf{z}_i^T W_{y_i}^{(k)}}{\|\mathbf{z}_i\|_2\|W_{y_i}^{(k)}\|_2}.
\end{equation}
Sau đó, mô hình áp dụng cùng một biên góc tương tự như ArcFace. Cơ chế đa tâm này rất phù hợp với đặc thù bộ dữ liệu S3, nơi các hình ảnh của một loài gỗ có thể tạo ra phân bố đa đỉnh trong không gian đặc trưng do sự khác biệt giữa gỗ lõi và gỗ giác, tuổi cây hoặc vị trí cắt. Thay vì ép buộc tất cả các biến thể tự nhiên này vào một trọng tâm duy nhất, SubCenter ArcFace cho phép mô hình tự động học các cụm phân bố cấu trúc bề mặt riêng biệt trong cùng một loài gỗ.

\subsubsection{Hàm loss SoftTriple}
Tương tự, phương pháp SoftTriple cũng sử dụng nhiều trọng tâm đại diện cho mỗi lớp nhưng áp dụng cơ chế gán mềm (soft assignment) thay vì lấy giá trị cực đại:
\begin{equation}
S_{i,c}=\sum_{k=1}^{K}\frac{e^{\mathbf{z}_i^T w_c^{(k)}/\gamma}}{\sum_{t=1}^{K}e^{\mathbf{z}_i^T w_c^{(t)}/\gamma}}\mathbf{z}_i^T w_c^{(k)}.
\end{equation}
Loss phân loại có margin:
\begin{equation}
L_{\text{SoftTriple}}=-\log\frac{e^{\lambda(S_{i,y_i}-\delta)}}{e^{\lambda(S_{i,y_i}-\delta)}+\sum_{c\ne y_i}e^{\lambda S_{i,c}}}.
\end{equation}
Cơ chế gán mềm giúp dòng đạo hàm cập nhật đồng thời cho nhiều trọng tâm trong mỗi bước huấn luyện, hạn chế hiện tượng trọng tâm bị cô lập hoặc không được cập nhật (dead center). Đối với hình ảnh gỗ, SoftTriple mang lại ý nghĩa thực tế lớn do các phân nhóm vân gỗ thường không phân tách một cách rõ rệt mà có sự chuyển tiếp liên tục giữa các vùng mô trên mặt cắt.

\subsubsection{Hàm loss tương phản có giám sát (SupCon)}
Hàm loss tương phản có giám sát (Supervised Contrastive Loss - SupCon) mở rộng phương pháp SimCLR sang học có giám sát bằng cách coi tất cả các mẫu thuộc cùng một loài trong lô dữ liệu là các mẫu tích cực của nhau:
\begin{equation}
L_{\text{SupCon}}=\sum_{i\in I}\frac{-1}{|P(i)|}\sum_{p\in P(i)}\log\frac{e^{\mathbf{z}_i^T\mathbf{z}_p/\tau}}{\sum_{a\in A(i)}e^{\mathbf{z}_i^T\mathbf{z}_a/\tau}}.
\end{equation}
SupCon tối ưu hóa việc sắp đặt đặc trưng trên toàn bộ lô dữ liệu và thường tạo ra không gian nhúng có khả năng chuyển giao tốt. Trong bài toán S3, phương pháp này giúp kéo hình ảnh của cùng một loài lại gần nhau bất kể những biến thiên do góc chụp, nhưng vẫn tiềm ẩn nguy cơ đẩy các mẫu có độ tương đồng cao sang hai phía quá mạnh (false-negative repulsion) when two species under the same genus are extremely similar.

\subsection{Các công thức hàm loss trong học tự giám sát}
Trong học biểu diễn tự giám sát (Self-Supervised Learning - SSL), mô hình thực hiện tối ưu hóa đặc trưng thông qua các biến thể tăng cường ngẫu nhiên của hình ảnh mà không cần sử dụng nhãn lớp giám sát. Phần này trình bày các hàm loss của bốn phương pháp học tự giám sát được khảo sát trong nghiên cứu.

\subsubsection{Hàm loss SimCLR}
Phương pháp SimCLR (Simple Framework for Contrastive Learning of Visual Representations) tối ưu hóa hàm loss tương phản InfoNCE trên các cặp ảnh tăng cường. Với kích thước lô dữ liệu là $N$, mỗi hình ảnh được tăng cường ngẫu nhiên thành 2 phiên bản (views), tạo thành một lô gồm $2N$ mẫu dữ liệu. Cặp mẫu tích cực (positive pair) là $(i,j)$ gồm hai phiên bản tăng cường khác nhau của cùng một hình ảnh gốc. Hàm loss cho mỗi cặp mẫu tích cực $(i,j)$ được phát biểu:
\begin{equation}
\begin{aligned}
\ell(i, j) = -\log \frac{e^{\mathbf{z}_i^T\mathbf{z}_j/\tau}}{\Omega_i},
\end{aligned}
\end{equation}
where:
\begin{equation}
\begin{aligned}
\Omega_i = \sum_{k=1}^{2N} \mathbb{I}_{[k \ne i]} e^{\mathbf{z}_i^T\mathbf{z}_k/\tau}.
\end{aligned}
\end{equation}
trong đó $\tau$ là siêu tham số nhiệt độ (temperature). Hàm loss tổng thể được tính bằng trung bình cộng của toàn bộ các cặp mẫu tích cực trong lô dữ liệu:
\begin{equation}
L_{\text{SimCLR}} = \frac{1}{2N} \sum_{k=1}^{N} \left[ \ell(2k-1, 2k) + \ell(2k, 2k-1) \right].
\end{equation}

\subsubsection{Hàm loss BYOL}
Phương pháp BYOL (Bootstrap Your Own Latent) sử dụng hai mạng nơ-ron song song: mạng học trực tiếp (online network) với tham số $\theta$ và mạng mục tiêu (target network) với tham số $\xi$. Mạng học trực tiếp có nhiệm vụ dự đoán biểu diễn của mạng mục tiêu từ một phiên bản ảnh khác. BYOL tối ưu hóa sai số bình phương trung bình (MSE) giữa vector dự đoán đã chuẩn hóa của mạng trực tiếp và vector đặc trưng của mạng mục tiêu:
\begin{equation}
L_{\text{MSE}}(q_\theta(z_\theta), z'_\xi) = \left\| \bar{q}_\theta - \bar{z}'_\xi \right\|_2^2 = 2 - 2 \bar{q}_\theta^T \bar{z}'_\xi,
\end{equation}
where:
\begin{equation}
\bar{q}_\theta = \frac{q_\theta(z_\theta)}{\|q_\theta(z_\theta)\|_2},\quad \bar{z}'_\xi = \frac{z'_\xi}{\|z'_\xi\|_2}.
\end{equation}
Việc ngăn chặn hiện tượng sụp đổ biểu diễn đặc trưng (representation collapse) được đảm bảo bằng cơ chế dừng cập nhật đạo hàm (\textit{stop-gradient}) đặt trên mạng mục tiêu:
\begin{equation}
L_{\text{BYOL}} = L_{\text{MSE}}(q_\theta(z_\theta), \text{stop\_gradient}(z'_\xi)).
\end{equation}
Tham số mạng mục tiêu $\xi$ được cập nhật chậm theo cơ chế trung bình trượt lũy thừa (EMA) từ $\theta$:
\begin{equation}
\xi \leftarrow m \xi + (1-m) \theta,
\end{equation}
trong đó $m \in [0, 1]$ là hệ số momentum.

\subsubsection{Hàm loss SimSiam}
Phương pháp SimSiam (Simple Siamese) đơn giản hóa kiến trúc BYOL bằng cách loại bỏ hoàn toàn mạng mục tiêu động lượng (momentum target network, tức thiết lập $\xi = \theta$). SimSiam chứng minh rằng chỉ cần sử dụng cơ chế dừng cập nhật đạo hàm kết hợp với kiến trúc không đối xứng (bộ dự đoán predictor) là đủ để ngăn chặn sự sụp đổ biểu diễn đặc trưng. Hàm loss đối xứng được định nghĩa như sau:
\begin{equation}
L_{\text{SimSiam}} = \frac{1}{2} \mathcal{D}\big(p_1, \text{sg}(z_2)\big) + \frac{1}{2} \mathcal{D}\big(p_2, \text{sg}(z_1)\big),
\end{equation}
trong đó $p_1 = q_\theta(z_1)$ và $p_2 = q_\theta(z_2)$ là đầu ra của bộ dự đoán cho hai phiên bản ảnh, còn $z_1, z_2$ là đầu ra của bộ đầu chiếu. $\text{sg}$ biểu thị toán tử dừng đạo hàm (stop-gradient). $\mathcal{D}$ là độ tương đồng cosine âm:
\begin{equation}
\mathcal{D}(p, z) = -\frac{p^T z}{\|p\|_2 \|z\|_2}.
\end{equation}

\subsubsection{Hàm loss Barlow Twins}
Phương pháp Barlow Twins áp dụng nguyên lý giảm thiểu thông tin dư thừa (redundancy reduction) lên ma trận tương quan chéo thu được từ hai phiên bản ảnh tăng cường. Gọi $C$ là ma trận tương quan chéo kích thước $D \times D$ tính toán dọc theo chiều lô dữ liệu $B$ giữa hai tập đặc trưng chiếu $\mathbf{z}^A$ và $\mathbf{z}^B$:
\begin{equation}
C_{ij} = \frac{\sum_{b=1}^{B} z_{b,i}^A z_{b,j}^B}{\sqrt{\sum_{b=1}^{B} (z_{b,i}^A)^2} \sqrt{\sum_{b=1}^{B} (z_{b,j}^B)^2}},
\end{equation}
trong đó $i, j$ đại diện cho các chiều của không gian đặc trưng nhúng. Hàm mất mát Barlow Twins được định nghĩa là:
\begin{equation}
L_{\text{BarlowTwins}} = \sum_{i=1}^{D} (1 - C_{ii})^2 + \lambda \sum_{i=1}^{D} \sum_{j \neq i} C_{ij}^2,
\end{equation}
trong đó số hạng thứ nhất (invariance term) ép các hệ số tương quan chéo trên đường chéo chính tiến về 1 (tối đa hóa độ tương đồng giữa hai phiên bản ảnh), số hạng thứ hai (redundancy reduction term) ép các hệ số tương quan chéo ngoài đường chéo chính tiến về 0 với hệ số phạt $\lambda > 0$ (giảm thiểu sự trùng lặp thông tin giữa các chiều đặc trưng nhúng).

\section{Thiết lập thực nghiệm}

\subsection{Chi tiết bộ dữ liệu S3}
Bộ dữ liệu hình ảnh gỗ S3 chứa hình ảnh mặt cắt ngang ở mức vĩ mô (macro) của 19 loài gỗ thuộc 6 chi khác nhau: \textit{Afzelia}, \textit{Dalbergia}, \textit{Guibourtia}, \textit{Peltogyne}, \textit{Pterocarpus} và \textit{Sindora}. Bộ dữ liệu có cấu trúc phân cấp tự nhiên: mỗi loài gỗ tương ứng với một lớp (class); bên trong mỗi loài chứa nhiều thư mục con đại diện cho một mẫu vật lý riêng biệt. Việc giữ nguyên vẹn toàn bộ hình ảnh của cùng một thư mục con vào cùng một tập dữ liệu là bắt buộc để tránh rò rỉ đặc trưng lối tắt từ phiên chụp. Các mẫu gỗ chỉ được định danh ở cấp chi được loại bỏ khỏi quy trình huấn luyện để đảm bảo tính nhất quán của nhãn.

\subsection{Tăng cường dữ liệu}
Quy trình tăng cường dữ liệu được thiết kế nhằm bảo toàn các cấu trúc giải phẫu sinh học (vân gỗ, mạch gỗ) nhưng nâng cao khả năng chống chịu trước các biến thiên thực tế: xoay ngẫu nhiên, cắt và thay đổi kích thước (RandomResizedCrop), biến đổi màu sắc nhẹ (ColorJitter), chuyển đổi ảnh xám, làm mờ Gaussian (GaussianBlur), điều chỉnh độ sắc nét và xóa ngẫu nhiên (RandomErasing). Các phép biến đổi này mô phỏng sự khác biệt về hướng cắt, độ phóng đại ống kính, độ sáng và các yếu tố nhiễu bề mặt.

\subsection{Giao thức thực nghiệm đối chiếu}
Nghiên cứu tiến hành đánh giá thực nghiệm đối chiếu 14 phương pháp học khoảng cách: Triplet cơ bản, Triplet bán khó, Triplet biên mềm, Angular, Multi-Similarity, Contrastive cơ bản, Contrastive khai thác cặp khó, Lifted Structured, Circle, Proxy Anchor, ArcFace, SubCenter ArcFace, SoftTriple và tương phản có giám sát (SupCon). Nhóm tự giám sát bao gồm bốn mô hình: SimCLR, BYOL, SimSiam và Barlow Twins. Tất cả các phương pháp này đều được huấn luyện và kiểm thử trên cùng một cấu hình chia dữ liệu thu được từ giao thức CEGS-Split để đảm bảo tính khách quan và loại trừ các sai số do phân bổ dữ liệu.

\subsection{Các chỉ số đánh giá hiệu năng không gian đặc trưng}
Chất lượng của không gian nhúng biểu diễn đặc trưng được kiểm chứng qua hai nhóm chỉ số đo lường chính. Nhóm truy xuất (retrieval) bao gồm: Recall@k, Precision@k, mAP và AUC. Đối với từng chỉ số đo lường theo từng lớp $x_i$, báo cáo sử dụng cả giá trị trung bình cộng (macro average) và trung bình điều hòa (harmonic mean) kèm theo một ngưỡng làm mịn $\epsilon=0.01$ để tránh chia cho 0:
\begin{equation}
\text{HM}=\frac{N}{\sum_{i=1}^{N}\frac{1}{x_i+\epsilon}},\quad \epsilon=0.01.
\end{equation}
Chỉ số trung bình điều hòa (Harmonic mean) giúp phản ánh chính xác hiệu năng ở các lớp gỗ yếu, tránh tình trạng mô hình đạt điểm số trung bình cộng (macro average) cao nhờ các loài dễ phân biệt nhưng lại hoàn toàn thất bại ở các loài gỗ cùng chi có độ tương đồng hình thái cao (fine-grained).

Nhóm clustering gồm Silhouette Score, Davies-Bouldin Index, Calinski-Harabasz Index, Dunn Index và Normalized Mutual Information. Với một mẫu $i$, Silhouette được tính:
\begin{equation}
s(i)=\frac{b(i)-a(i)}{\max(a(i),b(i))},
\end{equation}
trong đó $a(i)$ là khoảng cách trung bình đến các mẫu cùng cụm và $b(i)$ là khoảng cách trung bình nhỏ nhất đến cụm khác. Davies-Bouldin Index là
\begin{equation}
\text{DBI}=\frac{1}{C}\sum_{i=1}^{C}\max_{j\ne i}\frac{\sigma_i+\sigma_j}{d(\mu_i,\mu_j)}.
\end{equation}
Dunn Index được định nghĩa:
\begin{equation}
\text{DI}=\frac{\min_{i\ne j}\delta(C_i,C_j)}{\max_k\Delta(C_k)}.
\end{equation}
NMI đo mức tương đồng thông tin giữa nhãn cụm dự đoán và nhãn loài thật:
\begin{equation}
\text{NMI}(Y,C)=\frac{2I(Y;C)}{H(Y)+H(C)}.
\end{equation}

\subsection{Siêu tham số thực thi hệ thống}
Mô hình được tối ưu hóa bằng thuật toán AdamW với tốc độ học (learning rate) $10^{-4}$ và hệ số suy giảm trọng số (weight decay) $10^{-2}$. Bộ lập lịch suy giảm tốc độ học theo hàm cosin (Cosine Annealing scheduler) được sử dụng trong tối đa 100 chu kỳ huấn luyện (epoch). Cơ chế dừng sớm (early stopping) giám sát chỉ số harmonic mAP trên tập kiểm định với ngưỡng chờ (patience) là 30 chu kỳ. Số chiều đặc trưng nhúng (embedding dimension) mặc định là 256, tỷ lệ đóng băng các lớp mạng (freeze ratio) mặc định là $90\%$, và mạng cơ sở (backbone) chính được sử dụng là ConvNeXt-Tiny. Các thực nghiệm loại bỏ thành phần (ablation study) được thiết kế chi tiết để đánh giá ảnh hưởng của cấu trúc mạng cơ sở, số chiều đặc trưng nhúng, tỷ lệ đóng băng lớp mạng, phương pháp phân chia dữ liệu và hàm loss.

\section{Kết quả và Thảo luận}

\subsection{Đánh giá định lượng hiệu năng}
Bảng \ref{tab:main_results} trình bày khung báo cáo kết quả thực nghiệm chính của các phương pháp học biểu diễn đặc trưng trên tập kiểm thử. Khi cập nhật số liệu thực nghiệm thực tế, mỗi phương pháp cần được báo cáo đồng thời theo hai chỉ số trung bình cộng (macro average) và trung bình điều hòa (harmonic mean) để phản ánh rõ nét cả hiệu năng tổng thể lẫn độ công bằng, ổn định của mô hình đối với các loài gỗ hiếm hoặc có độ khó phân loại cao.

\begin{table*}[htbp]
\centering
\caption{Performance Comparison of Representation Learning Methods on the Test Set (CEGS-Split)}
\label{tab:main_results}
\resizebox{\textwidth}{!}{%
\begin{tabular}{lccccccccccc}
\toprule
\multirow{2}{*}{\textbf{Method}} & \multicolumn{2}{c}{\textbf{Recall@1 (\%)}} & \multicolumn{2}{c}{\textbf{Precision@1 (\%)}} & \multicolumn{2}{c}{\textbf{mAP (\%)}} & \multicolumn{5}{c}{\textbf{Clustering Metrics}} \\
\cmidrule(lr){2-3} \cmidrule(lr){4-5} \cmidrule(lr){6-7} \cmidrule(lr){8-12}
& Macro & Harmonic & Macro & Harmonic & Macro & Harmonic & Silhouette $\uparrow$ & DBI $\downarrow$ & CHI $\uparrow$ & Dunn $\uparrow$ & NMI $\uparrow$ \\
\midrule
\multicolumn{12}{l}{\textit{Classification Base}} \\
Focal Loss Baseline & - & - & - & - & - & - & - & - & - & - & - \\
\midrule
\multicolumn{12}{l}{\textit{Pairwise \& Triplet Metric Learning}} \\
Contrastive Loss & - & - & - & - & - & - & - & - & - & - & - \\
Hard Contrastive Loss & - & - & - & - & - & - & - & - & - & - & - \\
Triplet Loss & - & - & - & - & - & - & - & - & - & - & - \\
Semi-hard Triplet Loss & - & - & - & - & - & - & - & - & - & - & - \\
Soft-Margin Triplet & - & - & - & - & - & - & - & - & - & - & - \\
Lifted Structured Loss & - & - & - & - & - & - & - & - & - & - & - \\
Multi-Similarity Loss & - & - & - & - & - & - & - & - & - & - & - \\
Circle Loss & - & - & - & - & - & - & - & - & - & - & - \\
\midrule
\multicolumn{12}{l}{\textit{Angular Margin and Proxy-based Learning}} \\
Angular Loss & - & - & - & - & - & - & - & - & - & - & - \\
Proxy Anchor & - & - & - & - & - & - & - & - & - & - & - \\
ArcFace & - & - & - & - & - & - & - & - & - & - & - \\
SubCenter ArcFace & - & - & - & - & - & - & - & - & - & - & - \\
SoftTriple & - & - & - & - & - & - & - & - & - & - & - \\
\midrule
\multicolumn{12}{l}{\textit{Self-Supervised Learning}} \\
SupCon & - & - & - & - & - & - & - & - & - & - & - \\
SimCLR & - & - & - & - & - & - & - & - & - & - & - \\
BYOL & - & - & - & - & - & - & - & - & - & - & - \\
SimSiam & - & - & - & - & - & - & - & - & - & - & - \\
Barlow Twins & - & - & - & - & - & - & - & - & - & - & - \\
\midrule
\textbf{Proposed CEGS-Split (Ours)} & \textbf{-} & \textbf{-} & \textbf{-} & \textbf{-} & \textbf{-} & \textbf{-} & \textbf{-} & \textbf{-} & \textbf{-} & \textbf{-} & \textbf{-} \\
\bottomrule
\end{tabular}%
}
\end{table*}

\subsection{Ablation on Data Splitting Impact}
Để định lượng tác động của rò rỉ dữ liệu, chúng tôi so sánh Random Split với CEGS-Split bằng một giao thức KNN trên embedding Swin-Large. Random Split được chạy trên ba seed $(42,123,456)$, trong khi CEGS-Split dùng cấu hình class-wise đã chọn. Các chỉ số gồm KNN test accuracy, F1-Macro, F1-Weighted, mean cosine similarity và max cosine similarity giữa train-test. Mean cosine similarity phản ánh mức gần trung bình giữa ảnh test và không gian train; max cosine similarity đặc biệt quan trọng vì nó phát hiện các test sample có láng giềng gần bất thường trong train, một dấu hiệu của near-duplicate hoặc rò rỉ ngữ nghĩa.

\begin{table}[H]
\centering
\caption{Data splitting ablation using KNN on Swin-Large embeddings}
\label{tab:split_ablation}
\resizebox{\columnwidth}{!}{%
\begin{tabular}{lccccc}
\toprule
\textbf{Split Method} & \textbf{KNN Acc.} & \textbf{F1-Macro} & \textbf{F1-Weighted} & \textbf{Mean CosSim} & \textbf{Max CosSim} \\
\midrule
Random & $0.9581 \pm 0.0155$ & $0.9493 \pm 0.0179$ & $0.9571 \pm 0.0159$ & $0.5730 \pm 0.0037$ & $0.9161 \pm 0.0015$ \\
CEGS-Split & $0.9087 $ & $0.8858 $ & $0.9062 $ & $0.5711 $ & $0.8999 $ \\
\bottomrule
\end{tabular}%
}
\end{table}

Kết quả thực nghiệm ở Bảng \ref{tab:split_ablation} cho thấy phương pháp phân chia ngẫu nhiên (Random Split) đạt độ chính xác KNN trung bình là $0.9581$, cao hơn giao thức CEGS-Split ($0.9087$) một khoảng $0.0494$ (tương đương mức đánh giá quá cao hiệu năng thực tế là $4.94\%$). Chỉ số F1-Macro cũng sụt giảm đáng kể từ $0.9493$ xuống $0.8858$ khi chuyển đổi sang giao thức cô lập cấp mẫu vật vật lý, chứng minh rằng sự sụt giảm hiệu năng không chỉ xảy ra ở các loài gỗ phổ biến mà còn ảnh hưởng trực tiếp đến độ chính xác và tính công bằng ở các loài gỗ hiếm. Đồng thời, Random Split ghi nhận độ tương đồng cosine trung bình cao hơn CEGS-Split $0.0019$ và độ tương đồng cosine lớn nhất cao hơn $0.0162$ ($0.9161$ so với $0.8999$). Điều này chỉ ra rằng tập kiểm thử của phép chia ngẫu nhiên có xu hướng nằm quá gần tập huấn luyện trong không gian đặc trưng biểu diễn, dẫn đến việc mô hình hoặc bộ phân loại KNN dễ dàng đạt điểm số cao nhờ học thuộc lòng các đặc trưng lối tắt từ các ảnh có cùng nguồn gốc tiêu bản gỗ. Do đó, giao thức CEGS-Split thiết lập một môi trường đánh giá khách quan, khó khăn hơn nhưng phản ánh chính xác khả năng ứng dụng thực tế trên các khúc gỗ hoàn toàn mới.

\subsection{Thử nghiệm loại bỏ thành phần trên các cấu phần mô hình}
Các thực nghiệm loại bỏ thành phần (ablation study) dự kiến bao gồm việc đánh giá sự thay đổi hiệu năng khi thay thế cấu trúc mạng trích xuất đặc trưng cơ sở (ResNet-50, EfficientNetV2-S, ConvNeXt-Tiny, Swin-Tiny và DINOv2-ViT-S/14); thay đổi số chiều không gian nhúng ($64, 128, 256, 512$); thay đổi tỷ lệ đóng băng lớp mạng ($0\%, 50\%, 75\%, 90\%, 100\%$); và so sánh chéo giữa các hàm loss cốt lõi như Triplet, ArcFace, Multi-Similarity và SupCon. Kết quả thực nghiệm nên được trình bày dưới dạng giá trị trung bình kèm độ lệch chuẩn (mean $\pm$ standard deviation) thu được trên nhiều lượt huấn luyện ngẫu nhiên khác nhau, kết hợp sử dụng các phép kiểm định thống kê cặp (như paired t-test hoặc Wilcoxon signed-rank test) để đảm bảo các cải tiến hiệu năng có ý nghĩa thống kê thực sự chứ không phải do nhiễu thực nghiệm.

\section{Trực quan hóa không gian nhúng và Tính giải thích của mô hình}

\subsection{Trực quan hóa phân cụm bằng t-SNE}
Phân tích không gian đặc trưng bằng t-SNE tập trung làm rõ chi gỗ \textit{Afzelia}, bao gồm bốn loài có mức độ tương đồng giải phẫu học cao: \textit{A. africana}, \textit{A. bella}, \textit{A. pachyloba} và \textit{A. quanzensis}. Đây là một nhóm phân loại cực kỳ khó phân biệt do cấu trúc hình thái mạch gỗ và tia gỗ của chúng gần như giống hệt nhau. Đối với mô hình phân loại truyền thống hoặc các hàm loss học khoảng cách thông thường, đặc trưng nhúng của bốn loài gỗ này thường bị chồng lấn rất mạnh, tạo nên các vùng phân bố mơ hồ và dễ gây nhầm lẫn. Sau khi tối ưu hóa không gian nhúng bằng các phương pháp học khoảng cách sâu đề xuất, kết quả kỳ vọng là các mẫu gỗ vẫn duy trì phân nhóm chung cấp chi nhưng tự động tách thành bốn cụm loài riêng biệt rõ rệt, chứng minh khả năng trích xuất các đặc trưng mịn siêu vi mô của mô hình.

\subsection{Bản đồ chú ý giải thích bằng Grad-CAM}
Các phương pháp trực quan hóa bản đồ chú ý như Grad-CAM và Finer-CAM được áp dụng để phân tích các vùng điểm ảnh mà mô hình tập trung quan sát khi đưa ra quyết định phân loại loài gỗ. Trước khi thực hiện tinh chỉnh (fine-tuning), mạng cơ sở thường chỉ chú ý vào các thông tin thô như màu sắc tổng thể, vùng nền ảnh hoặc các chi tiết tạp chất trên bề mặt lát cắt. Sau khi tối ưu hóa bằng các hàm loss học khoảng cách sâu, bản đồ kích hoạt chú ý của mô hình dịch chuyển rõ rệt về các cấu trúc giải phẫu sinh học quan trọng của thớ gỗ (như hình dạng lỗ mạch, mật độ phân bố nhu mô quanh mạch, hướng đi của tia gỗ và cấu trúc sợi gỗ). Việc tập trung sự chú ý vào các marker giải phẫu này giúp khẳng định mô hình đã học được các tri thức khoa học thực sự có tính diễn giải sinh học, thay vì dựa vào các đặc trưng lối tắt phi bản chất để đạt độ chính xác cao.

\section{Kết luận và Định hướng tương lai}
Nghiên cứu này đề xuất một khung phương pháp luận toàn diện cho bài toán nhận diện loài gỗ ở cấp độ chi tiết, kết hợp chặt chẽ giữa giao thức chia dữ liệu chống rò rỉ và các thuật toán học khoảng cách sâu. Giao thức CEGS-Split chuyển đổi quy trình đánh giá truyền thống từ phân chia ngẫu nhiên cấp ảnh sang phân chia cô lập ở cấp mẫu vật vật lý độc lập, đảm bảo kết quả đo lường phản ánh chính xác khả năng tổng quát hóa của hệ thống trên các mẫu gỗ thực tế.

Trong tương lai, nghiên cứu sẽ được mở rộng theo hướng đánh giá chéo giữa các bộ dữ liệu khác nhau (cross-dataset validation) với các kho dữ liệu gỗ quốc tế như UFPR, bổ dung kiểm định thống kê trên nhiều giá trị seed ngẫu nhiên khác nhau, đánh giá độ bền bỉ (robustness) trước các thay đổi về thiết bị thu nhận ảnh và phát triển các giao thức nhận diện tập mở (open-set) để phát hiện các loài gỗ mới chưa từng xuất hiện trong tập huấn luyện. Một định hướng quan trọng khác là tích hợp các tri thức mô tả đặc điểm giải phẫu gỗ từ chuyên gia vào quá trình tối ưu hóa mạng nơ-ron, nhằm gia tăng tính diễn giải khoa học và độ tin cậy pháp lý trong các ứng dụng thực tế của pháp y lâm nghiệp.

\begin{thebibliography}{00}
\bibitem{b1} CITES, ``Convention on International Trade in Endangered Species of Wild Fauna and Flora,'' 1973.
\bibitem{b2} A. L. Hafemann, L. S. Oliveira, and P. Cavalin, ``Forest species identification using deep convolutional neural networks,'' in \textit{Proceedings of the International Conference on Pattern Recognition (ICPR)}, 2014, pp. 1107--1112.
\bibitem{b3} J. Deng, J. Guo, N. Xue, and S. Zafeiriou, ``ArcFace: Additive Angular Margin Loss for Deep Face Recognition,'' in \textit{Proceedings of the IEEE/CVF Conference on Computer Vision and Pattern Recognition (CVPR)}, 2019, pp. 4690--4699.
\bibitem{b4} Y. Wang, Q. Yao, and J. T. Kwok, ``Multi-similarity loss with general pair weighting for deep metric learning,'' in \textit{Proceedings of the IEEE/CVF Conference on Computer Vision and Pattern Recognition (CVPR)}, 2019, pp. 5022--5030.
\bibitem{b5} P. Khosla et al., ``Supervised contrastive learning,'' \textit{Advances in Neural Information Processing Systems (NeurIPS)}, vol. 33, pp. 18661--18673, 2020.
\bibitem{b6} J. Zbontar et al., ``Barlow Twins: Self-supervised learning via redundancy reduction,'' in \textit{Proceedings of the International Conference on Machine Learning (ICML)}, 2021.
\bibitem{b7} Q. Qian et al., ``SoftTriple Loss: Deep Metric Learning Without Triplet Sampling,'' in \textit{Proceedings of the IEEE/CVF International Conference on Computer Vision (ICCV)}, 2019.
\bibitem{b8} Y. Sun et al., ``Circle Loss: A unified perspective of pair similarity optimization,'' in \textit{CVPR}, 2020.
\end{thebibliography}

\end{document}
